% SHARD-ID: AQM-58-PRODUCT-RINGS-REAL-LINE-COTANGENT
% SHARD-TITLE: Product Rings and the Formal Real Affine Line
% SHARD-SUMMARY: Proves that \(k^n\) has trivial K\"ahler differentials and hence no de Rham CSS dynamics.
% SHARD-SUMMARY: Computes \(\Omega^1_{\mathbb R[x]/\mathbb R}\), cotangent fibres, and tangent vectors for the real affine line.
% SHARD-SUMMARY: Separates the algebraic cotangent phase space from the analytic Schr\"odinger quantisation choices.
% SHARD-KEYWORDS: product ring, real line, cotangent, tangent, Kahler differentials, de Rham CSS, Schrodinger
\section{Product Rings and the Formal Real Affine Line}
\label{sec:product-rings-real-line-cotangent}
\subsection{Status and Source Anchors}
This shard answers two follow-up questions about convention
\texttt{CONVENTIONS.md} (ao).  K\"ahler differentials, polynomial
differentials, and the de Rham complex use the Stacks algebra snapshot
\path{references/algebraic_geometry/stacks_project_algebra.tex}, lines
33793--33872, 34310--34340, and 34388--34524.  Tangent space by dual numbers
and cotangent-tangent duality use
\path{references/algebraic_geometry/stacks_project_varieties.tex}, lines
2915--3019.  The product-ring Weyl layer is AQM-40 and convention (ad).  The
real Weyl/Schr\"odinger representation statement is the finite-dimensional
CCR layer of AQM-19, sourced there from Derezinski and Bekka.
\subsection{Part I: Product Rings Have Trivial Derivative Dynamics}
Let \(k\) be a field and
\[
  A=k^n=k e_1\oplus\cdots\oplus k e_n,
\]
where \(e_i\) is the idempotent with \(1\) in the \(i\)-th factor and \(0\)
elsewhere.  Let \(M\) be any \(A\)-module and let \(D:A\to M\) be a
\(k\)-derivation.  We first prove \(D(e_i)=0\) without assuming anything
about the characteristic.  Since \(e_i^2=e_i\),
\[
  D(e_i)=D(e_i^2)=e_iD(e_i)+e_iD(e_i).
\]
Multiplying by \(e_i\) gives
\[
  e_iD(e_i)=2e_iD(e_i),
\]
so \(e_iD(e_i)=0\).  Multiplying the displayed derivation equation by
\(1-e_i\) gives
\[
  (1-e_i)D(e_i)=0.
\]
Adding the two pieces gives \(D(e_i)=0\).

Every element of \(A\) is \(a=\sum_i a_i e_i\) with \(a_i\in k\).  Because
\(D\) is a \(k\)-derivation, \(D(a_i)=0\), and hence
\[
  D(a)=\sum_i a_iD(e_i)=0.
\]
Thus \(\operatorname{Der}_k(A,M)=0\) for every \(M\).  By the universal
property,
\[
  \Omega^1_{A/k}=0.
\]
Consequently
\[
  A\longrightarrow 0\longrightarrow0
\]
is the de Rham complex, and the de Rham CSS test has
\[
  C^1=0,\qquad R_X=0,\qquad R_Z=0.
\]
So \(k^n\) has trivial derivative dynamics in exactly the same sense as the
single field \(k\).  This does not erase the kinematics.  By AQM-40,
\[
  \Spec(k^n)=\{x_1,\ldots,x_n\},\qquad \kappa(x_i)=k,
\]
and the residue-Weyl layer is the ambient \(n\)-qudit system
\[
  \mathcal H=\ell^2(k)^{\otimes n},\qquad
  \mathcal A=\bigotimes_{i=1}^n\operatorname{End}_{\mathbb C}\ell^2(k).
\]
The conclusion is therefore sharp: the product ring supplies positions and
the ambient Weyl/Pauli kinematics, but the K\"ahler derivative proposal
selects no stabilizer and no Hamiltonian.
\subsection{Part II: Cotangent Data on \texorpdfstring{\(\operatorname{Spec}\mathbb R[x]\)}{Spec R[x]}}
Let
\[
  A=\mathbb R[x],\qquad X=\Spec A.
\]
The polynomial differential source gives
\[
  \Omega^1_{A/\mathbb R}=A\,dx=\mathbb R[x]\,dx.
\]
The universal derivation is the formal derivative:
\[
  d(f(x))=f'(x)\,dx.
\]
The algebraic de Rham complex begins
\[
  \mathbb R[x]\xrightarrow{d_0}\mathbb R[x]\,dx\longrightarrow0.
\]
Since every polynomial \(g(x)\) has the polynomial antiderivative
\[
  \int g(x)\,dx
\]
over \(\mathbb R\), the map \(d_0\) is onto \(\mathbb R[x]\,dx\).  Its kernel
is the constants.  Hence
\[
  H^0_{\rm dR}(\mathbb R[x]/\mathbb R)=\mathbb R,\qquad
  H^1_{\rm dR}(\mathbb R[x]/\mathbb R)=0.
\]
This is only a formal algebraic calculation; it is not an analytic statement
about domains of differential operators.

Now compute fibres.  At the real closed point
\[
  a\in\mathbb R,\qquad \mathfrak m_a=(x-a),
\]
the residue field is \(\kappa(\mathfrak m_a)=\mathbb R\).  The cotangent
fibre is
\[
  T^*_{X/\mathbb R,a}
  =
  \Omega^1_{A/\mathbb R}\otimes_A\mathbb R
  =
  \mathbb R\,dx_a.
\]
The tangent space is the dual line
\[
  T_{X/\mathbb R,a}
  =
  \operatorname{Hom}_{\mathbb R}(\mathbb R\,dx_a,\mathbb R)
  \cong\mathbb R.
\]
The tangent vector \(v\) with \(v(dx_a)=1\) corresponds to the derivation
\[
  D_v(f)=f'(a).
\]
More generally \(c\,v\) corresponds to \(D_{c v}(f)=c f'(a)\).

There are also non-real closed points.  If
\[
  \pi(x)=x^2+b x+c
\]
is irreducible over \(\mathbb R\), then
\[
  \kappa(\pi)=\mathbb R[x]/(\pi)\cong\mathbb C,
\]
and
\[
  T^*_{X/\mathbb R,\pi}\cong \mathbb C\,dx_\pi,\qquad
  T_{X/\mathbb R,\pi}\cong \operatorname{Hom}_{\mathbb C}(\mathbb C\,dx_\pi,
  \mathbb C).
\]
The generic point has residue field \(\mathbb R(x)\), so its cotangent fibre is
\[
  T^*_{X/\mathbb R,\eta}\cong\mathbb R(x)\,dx.
\]
\subsection{Predicted Quantisation}
The cotangent construction says that the algebraic phase space of the real
affine line is the total cotangent bundle
\[
  T^*X=\Spec\operatorname{Sym}_A(\Omega^1_{A/\mathbb R})
  \cong \Spec\mathbb R[x,p].
\]
Here \(x\) is the position coordinate and \(p\) is the fibre-linear momentum
coordinate dual to \(dx\).  Its canonical algebraic symplectic form is
\[
  dp\wedge dx.
\]
Thus the cotangent-Weyl prediction is the usual one-dimensional CCR phase
space with labels \((q,p)\in\mathbb R^2\), symplectic form
\[
  \Omega((q,p),(q',p'))=p q'-p'q.
\]
By the AQM-19 finite-mode CCR convention, the irreducible regular
representation is the Schr\"odinger representation on \(L^2(\mathbb R)\):
position is multiplication by \(x\), and momentum is represented formally by
\[
  -i\,\frac{d}{dx}.
\]
The word ``formally'' is essential here.  The cotangent calculation supplies
the phase-space geometry and the CCR labels.  It does not by itself choose
self-adjoint domains, a Hamiltonian such as \(-d^2/dx^2+V(x)\), boundary
conditions, or a measure beyond the standard analytic Schr\"odinger choice.

The de Rham CSS proposal is different.  Applied literally to
\(\mathbb R[x]\), it is infinite-dimensional and outside the finite-qudit
scope.  Moreover \(d_0\) is surjective, so a naive finite-dimensional
truncation would tend to impose all \(C^1\) directions as \(X\)-checks.  The
physical particle-on-a-line quantisation therefore belongs to the
cotangent-Weyl branch, not to the finite de Rham CSS stabilizer test.
